\documentclass[10pt,a4paper]{amsart}
\usepackage[margin=19mm]{geometry}
\usepackage{amsmath,amssymb,mathtools}
\usepackage[scr=rsfs,cal=euler]{mathalfa}
\usepackage{xfrac}
\usepackage[unicode,hidelinks,bookmarks=false]{hyperref}
\newcommand{\ie}{i.e.\ }
\newcommand{\cf}{cf.\ }
\newcommand{\resp}{respectively\ }
\newcommand{\Subsection}{\subsection}
\providecommand{\nobreakdash}{\nobreak\hbox{-}}
\setlength{\emergencystretch}{2em}
\multlinegap0pt
\usepackage{bm}
\usepackage{bbm}
\usepackage{dsfont}
\usepackage{xspace}
\usepackage{cancel}
\usepackage{mathtools}
\usepackage{cleveref}
\usepackage{amsthm}
\makeatletter
\@ifundefined{theorem}{\newtheorem{theorem}{Theorem}}{}
\@ifundefined{lemma}{\newtheorem{lemma}[theorem]{Lemma}}{}
\makeatother
\newcommand{\Inv}{\textsc{inv}}
\newcommand{\blambda}{\bm{\lambda}}
\renewcommand{\u}{\mathsf{u}}
\title[An elaborate new proof of Cayley's formula]{An elaborate new proof of Cayley's formula}
\author{Esther Banaian, Anh Trong Nam Hoang, Elizabeth Kelley, Weston Miller, Jason Stack, Carolyn Stephen, Nathan Williams}
\date{Source: arXiv:2402.07798v1 (CC BY 4.0)}
\begin{document}
\maketitle

\noindent\emph{Source and licence.} An elaborate new proof of Cayley's formula. Version arXiv:2402.07798v1 (CC BY 4.0), distributed under the Creative Commons Attribution 4.0 International licence, as recorded in the arXiv licence metadata for this exact version and shown on the abstract page https://arxiv.org/abs/2402.07798v1. Full rights notice: licenses/arxiv-2402.07798-CC-BY-4.0.txt.\par\smallskip

\noindent\textit{Editorial scope.} Counted unit: the Lemma whose source label is \texttt{smallMoves} in the article (source0.tex lines 1591--1594, source label \texttt{smallMoves}), delivered together with the article's own complete original proof of it (source0.tex lines 1595--1632, one \texttt{proof} environment of 38 lines). No mathematics is written, repaired or re-derived: statement and proof are the article's own text line for line. Required context retained uncounted: maxDist (lines 1572--1579). Every definition, display and label of those lines is retained unchanged. Omitted from this package and not redistributed: 1--1571, 1580--1590, 1633--2237. Nothing inside a retained excerpt is omitted or altered: every retained body is delivered exactly as the article's lines are written, so no omission receipt of any kind accompanies this package. Citations: the retained excerpts carry no citation command at all, so no bibliography entry of the article is transcribed and no reference is redistributed. Reference adaptation: none was needed and none was made; every reference inside the delivered unit resolves inside it, so no unresolved reference or citation is produced. Printed slips kept literal: no printed word, symbol, hypothesis or number of the counted statement or of its proof was altered. Layout and class: the article's own class is \texttt{amsart} with packages including geometry, enumerate, amsmath, amsfonts, amssymb, amsthm, mathtools, thmtools, xparse, array, colortbl, hyperref, url, hypcap; the wrapper is the lane's pinned \texttt{amsart} editorial wrapper at 10pt on A4, which restates the theorem-like environments the retained text needs and adds the article's own macro definitions that the retained text needs (\texttt{\textbackslash Inv}, \texttt{\textbackslash blambda}, \texttt{\textbackslash u}), with extra wrapper packages (bm, bbm, dsfont, xspace, cancel, mathtools, cleveref). The delivered numbering is the unit's own. Assets: no retained span contains a figure environment, a graphics-inclusion command, a PDF-page inclusion, a table or a TikZ picture; no image file is copied into this package, the package declares no asset dependency and no image file is redistributed with it. Licence: version arXiv:2402.07798v1 of this article is distributed under the Creative Commons Attribution 4.0 International licence, as recorded in the arXiv licence metadata for this exact version and shown on the abstract page https://arxiv.org/abs/2402.07798v1. A full rights notice is delivered at licenses/arxiv-2402.07798-CC-BY-4.0.txt. Characters: every retained excerpt is pure ASCII, so the delivered engine, XeLaTeX with its default Latin Modern text font, reports no missing character.
\begin{lemma}
\label{maxDist}
    If $\u$ is an $e$-subword of $\blambda_n$, then for all integers $k$ and all $a, b$ with $0 \leq b - a < n(n-1)$,
    \[
        \big|\left(u_{a}\cdots u_{b}\right)(k) - k\big| \leq n-2,
    \]
    where the indices are taken modulo $n(n-1)$.
\end{lemma}

\begin{lemma}
\label{smallMoves}
    If $\u$ is an $e$-subword of $\blambda_n$ and $(\!(a,b)\!) \in \Inv(\u)$, then $|b - a| \leq n-1$.
\end{lemma}

\begin{proof}
    We need to show that
    \[
        \left|\left(u_1\cdots u_{j}\right)(j+1) - \left(u_1\cdots u_{j}\right)(j)\right| \leq n-1
    \]
    for all $j = 0,\dots,n(n-1)-1$. Suppose not. Since each inversion is an affine reflection, it is not possible to have
    \[
        |\left(u_1\cdots u_{j}\right)(j+1) - \left(u_1\cdots u_{j}\right)(j)| = n,
    \]
    so suppose that
    \[
        \left|\left(u_1\cdots u_{j}\right)(j+1) - \left(u_1\cdots u_{j}\right)(j)\right| \geq n+1.
    \]
    By \cref{maxDist} there are two cases to consider:
    \begin{enumerate}
        \item $\left(u_1\cdots u_{j}\right)(j+1) \leq j < j + 1 \leq \left(u_1\cdots u_{j}\right)(j)$
        \item $\left(u_1\cdots u_{j}\right)(j) \leq j < j + 1 \leq \left(u_1\cdots u_{j}\right)(j+1)$
    \end{enumerate}
    In the first case, there must exist an index $1 < \ell < j$ such that $u_\ell = s_{\ell - 1}$ and
    \[
        \left(u_{\ell + 1} \cdots u_j\right)(j) = \ell-1 \quad \text{and} \quad \left(u_{\ell + 1} \cdots u_j\right)(j+1) = \ell.
    \]
    Choose the smallest such $\ell$. Then
    \[
        \left|\left(u_{1}\cdots u_{\ell-1}\right)(\ell-1) - \left(u_{1}\cdots u_{\ell-1}(\ell)\right)\right| = \left|\left(u_1\cdots u_{j}(j+1)\right) - \left(u_1\cdots u_{j}\right)(j)\right| \geq n+1,
    \]
    and
    \[
        \left(u_1\cdots u_{\ell-1}\right)(\ell - 1) \leq \ell-1 < \ell \leq \left(u_1\cdots u_{j}\right)(\ell).
    \]
    So the first case reduces to the second.  Let
    \[
        c = j -  \left(u_1\cdots u_{j}\right)(j) \quad \text{and} \quad d = \left(u_1\cdots u_{j}\right)(j+1) - (j+1),
    \]
    so $c + d \geq n$. Now there must be indices $1 \leq i_d < i_{d-1} \cdots < i_1 \leq j$ such that $u_{i_\ell} = s_{j + \ell}$. Each of these indices must be on its own row. Moreover, $i_1$ cannot be on the same row as $j$.

    Now there must also be indices $j+1 \leq k_1 < \cdots < k_c \leq n(n-1)$ such that $u_{k_\ell} = s_{j - \ell}$. Again, each of these indices must be its own row. Moreover, $k_1$ cannot be on the same row as $j$.  It follows that we need to use at least $c + d + 1 \geq n + 1$ rows, which is a contradiction.
\end{proof}

\end{document}
